\textbf{Theorem 4} \ \textit{When no deadline is missed, a bound on the workload of a task $\tau_i$ in a generic interval $[a,b)$ can be computed considering a situation in which the carried-in job $J_i^\varepsilon$ starts executing at the beginning of the interval, with $a = d_i^\varepsilon - C_i$, and every other instance of $\tau_i$ is executed as soon as possible.}

\bigskip

\noindent
\textbf{Proof:} The situation is represented in Figure 3. Since a job $J_i^j$ can be ready only in $[r_i^j, d_i^j)$ and for at most $C_i$ time units, it is immediate to see that the depicted situation provides the highest amount of execution possible in interval $[a,b)$. Moving backwards the interval, the carry-in cannot increase, while the carry-out can only decrease. Instead, advancing the interval, the carry-in will decrease, while the carry-out can increase by at most the same amount. The situation is periodic. \hfill $\blacksquare$